
\section{Minimal Requirements for a Theory of Everything}
\label{sect:minimal-requirements}

We define a complete Theory of Everything ($\text{ToE}_{\text{complete}}$) as a formal framework that leaves no fundamental assumptions unexplained.
If a theory is structurally incomplete, it can be expressed as the complete theory missing a specific set of foundational elements:

\[
\text{ToE}_{\text{incomplete}} = \text{ToE}_{\text{complete}} \setminus \mathcal{A}
\]

where $\mathcal{A}$ denotes the set of fundamental assumptions left unexplained by the framework.
For a theory to truly be a theory of \textit{everything}, it must satisfy the condition of absolute explanatory self-sufficiency:

\[
\mathcal{A} = \emptyset
\]

By requiring $\mathcal{A}$ to be the empty set, the theory admits no external, unproven premises. 

